\section{Related Work}
\label{sec:related_work}

\subsection{Edge-Cloud IoT Task Offloading}

Mobile edge computing places compute and storage resources closer to users, devices, and sensors. Moving selected tasks away from distant cloud centers can reduce response time and improve local context awareness \cite{ref24,ref28}. The same architecture creates a scheduling problem. Edge resources are finite, and the network between devices, edge nodes, and cloud nodes is not stable. A task that is best served by the edge during light load may be better served by the cloud when the local queue grows.

Classical offloading studies model the communication-computation tradeoff directly. Binary offloading, energy-aware offloading, and joint radio-compute allocation have been studied with optimization methods and queue-aware formulations \cite{ref17,ref18,ref19}. These methods provide a strong analytical foundation. They also make the cost structure explicit: local execution depends on device or edge cycles, remote execution depends on transmission rate and server availability, and the final decision should respect latency and energy constraints. Their limitation in the present setting is not that the models are wrong, but that their assumptions can become brittle under bursty IoT traces, edge faults, cloud maintenance, and non-identical regional workloads.

Recent surveys and task-offloading studies show a clear movement from static optimization toward learning-based online scheduling \cite{ref12,ref28}. The reason is practical. In a real deployment, the scheduler may not know the exact future queue, channel, and task mix. A learning policy can adapt from observations. The open question is how to train such a policy when the data are naturally distributed across edge zones and when difficult tasks may be rejected because of deadline, battery, security, or resource constraints.

\subsection{Deep Reinforcement Learning for Offloading}

DRL is widely used for online offloading because it can map high-dimensional states to actions without requiring a closed-form solution for every operating condition \cite{ref03,ref07,ref12}. In these methods, the state often contains channel quality, queue length, task size, available compute, and energy indicators. The reward usually penalizes latency, energy consumption, or a weighted combination of both. This work keeps the broad MDP framing but gives explicit attention to rejection and SLA behavior, since a scheduler that reduces mean latency while increasing infeasible cases is not acceptable for QoS-sensitive IoT services.

DDQN fits binary or small discrete action spaces because it separates action selection from target evaluation, reducing the overestimation problem that can appear in standard DQN \cite{ref35,ref36}. Dueling Q-networks further separate the state value and the action advantage, which helps when a state is favorable or poor regardless of the chosen action \cite{ref41}. These properties match the Stage-1 offloading decision, where the action is binary and many states are dominated by task urgency or network quality.

Actor-critic methods and deterministic policy gradients are often used when the action is continuous or when resource allocation is optimized together with offloading \cite{ref38}. For that reason, the experimental comparison includes an FL-DDPG-style baseline. The difference in our design is deliberate: the proposed \fedddqn{} focuses on the binary offloading decision, and continuous-like CPU, memory, and bandwidth assignment is handled by a separate Stage-2 allocator. This split keeps the reinforcement-learning action space small and makes it easier to interpret the offloading decision itself.

\subsection{Federated Learning and Non-IID Edge Zones}

Federated learning trains a shared model from distributed clients without directly pooling raw local data \cite{ref30,ref32}. FedAvg remains the basic aggregation rule, while FedProx adds a proximal term to reduce client drift under heterogeneous data distributions \cite{ref30,ref31}. Edge-cloud IoT is a natural setting for such methods because local zones can differ in task types, radio quality, device classes, and failure patterns.

Federated reinforcement learning has recently been applied to task offloading and resource allocation in mobile edge and wireless IoT systems \cite{ref02,ref04,ref05,ref06}. These studies motivate the present work, but the evaluation focus here is different. Instead of treating federated learning mainly as a training architecture, the study asks whether a federated DDQN policy can make reliable edge-cloud decisions under deployment-inspired stress. The benchmark includes non-IID zones, workload phases, network outages, congestion windows, edge degradation, and cloud maintenance, so the federated setting affects the decision problem as well as the training layout.

\subsection{Joint Offloading and Resource Allocation}

Resource allocation is closely tied to offloading. A cloud or edge decision is incomplete if the scheduler ignores CPU cycles, memory demand, bandwidth, and queue pressure \cite{ref17,ref18,ref19}. Several studies optimize offloading and resource allocation jointly, often with energy, delay, or throughput objectives. Joint optimization can be powerful, but it can also produce a large action space and make it harder to isolate whether the offloading policy or the resource allocation rule caused a result.

The proposed method uses a two-stage formulation. The first stage learns the edge/cloud choice with a binary \fedddqn{} policy. The second stage applies a residual allocator only to edge-selected tasks. The allocator cannot hide a poor offloading policy by changing the destination. It can only improve proxy-target allocation fit for tasks already routed to the edge. That separation keeps the proposed offloading model clear while still adding a practical resource-management step.

\begin{table*}[!t]
\centering
\caption{Positioning of the proposed work relative to related edge-cloud offloading research.}
\label{tab:related_taxonomy}
\scriptsize
\begin{tabularx}{\textwidth}{@{}lXXX@{}}
\toprule
Category & Representative focus & Common limitation for this study & How this paper differs \\
\midrule
Optimization-based MEC offloading & Binary offloading, radio-compute allocation, energy and latency tradeoffs \cite{ref17,ref18,ref19} & Often relies on tractable model assumptions and may not adapt to burst/fault state changes. & Uses learned offloading under generated temporal bursts, congestion, queue carryover, and fault windows. \\
DRL offloading & DQN, DDQN, actor-critic, and online scheduling policies \cite{ref03,ref07,ref12} & Many evaluations emphasize average latency or energy without explicit rejection-aware analysis. & Uses reward and reporting that include SLA miss, rejection, edge pressure, and heavy-task routing behavior. \\
Federated edge intelligence & FedAvg, FedProx, and federated DRL across distributed clients \cite{ref02,ref04,ref05,ref06,ref30,ref31} & Federated structure can be shown without strong stress-test evidence or scenario-level diagnostics. & Trains across non-IID edge zones and reports scenario-wise, multi-seed, and ablation evidence. \\
Joint offloading and allocation & Coupled destination selection and resource assignment \cite{ref17,ref18,ref19} & Large joint action spaces can obscure whether gains come from routing or allocation. & Separates Stage-1 binary offloading from Stage-2 residual resource allocation for interpretability. \\
Benchmark and simulation studies & Synthetic tasks, simulated MEC nodes, and replayed state tables \cite{ref24,ref28} & Independent samples can understate contention and race-like load effects. & Uses \dataset{} with fixed schemas, time-correlated stress windows, and precomputed contention pressure. \\
\bottomrule
\end{tabularx}
\end{table*}


\subsection{Position of This Work}

The closest prior line of work combines DRL, MEC offloading, and federated training \cite{ref02,ref04,ref05,ref06}. The present work differs in four ways. First, it uses a fixed synthetic benchmark with burst, outage, queue, and maintenance effects instead of independent task samples. Second, it evaluates rejection as an outcome, not as an optional action that can be used to avoid difficult tasks. Third, it uses purged temporal train, validation, and test splits so the final comparison is not tuned on the test interval. Fourth, it adds a post-offloading allocator and evaluates that allocator separately from the Stage-1 policy. These choices give a controlled basis for comparing offloading behavior, QoS outcomes, and resource-allocation quality.
